%% kapa: exam=1 round=27 year=2016 period=1 form=A subject=economics subject_ko=경제학원론 seq=10 number=50
%% source: Q-Net 감정평가사 기출문제 (cst00309) artlSeq=5169402 file=extracted/1008/5169402/제27회 감정평가사 1차 문제/제27회 감정평가사 1교시문제(A형).pdf
%% conversion: kordoc 4.24.51 -> tools/gen1.py (auto); review=proofread-sonnet; flags=none
\begin{question}{50}
독점기업 甲은 두 시장 A, B에서 X재를 판매하고 있다. 생산에 있어서 甲의 한계비용은 0이다. 甲이 A, B에서 직면하는 수요함수는 각각 $Q_A = \mathrm{a}_1 - \mathrm{b}_1P_A$, $Q_B = \mathrm{a}_2 - \mathrm{b}_2P_B$이고, 甲이 각 시장에서 이윤극대화를 한 결과 두 시장의 가격이 같아지게 되는 $(\mathrm{a}_1, \mathrm{b}_1, \mathrm{a}_2, \mathrm{b}_2)$의 조건으로 옳은 것은? (단, $\mathrm{a}_1, \mathrm{b}_1, \mathrm{a}_2, \mathrm{b}_2$는 모두 양(+)의 상수이고, $Q_A$, $Q_B$는 각 시장에서 팔린 X재의 판매량이며, $P_A$, $P_B$는 각 시장에서 X재의 가격이다.)
\begin{choices}
  \item $\mathrm{a}_1\mathrm{a}_2 = \mathrm{b}_1\mathrm{b}_2$
  \item $\mathrm{a}_1\mathrm{b}_1 = \mathrm{a}_2\mathrm{b}_2$
  \item $\mathrm{a}_1\mathrm{b}_2 = \mathrm{a}_2\mathrm{b}_1$
  \item $\mathrm{a}_1+\mathrm{b}_1 = \mathrm{a}_2+\mathrm{b}_2$
  \item $\mathrm{a}_1+\mathrm{b}_2 = \mathrm{a}_2+\mathrm{b}_1$
\end{choices}
\end{question}
\answer{3}
